% ===== Korean script support for the Korean edition =====
% Included AFTER preamble.tex (see build_pdf.sh). ElegantBook has no Korean
% language option, so the class runs in `lang=en` (English theorem names,
% "Chapter"/"Contents" labels). Korean rendering is purely a font concern:
% the body/heading Latin fonts do not cover the Hangul blocks, and Noto
% Serif CJK KR does not cover Latin — so we keep the book's default
% Latin/CJK/symbol fonts and route ONLY Hangul codepoints to Noto Serif CJK KR.
%
% ucharclasses inserts the begin/end code at every entry/exit of a Hangul run.
% Wrapping in \begingroup…\endgroup means the exact surrounding style (bold in
% headings, monospace in code) is restored for whatever follows a Hangul run.
\usepackage{ucharclasses}

% Hangul Syllables (AC00–D7A3)
\IfFontExistsTF{Noto Serif CJK KR}{%
  \newfontfamily\hangulfont{Noto Serif CJK KR}%
  \setTransitionsFor{HangulSyllables}{\begingroup\hangulfont}{\endgroup}%
}{%
  \PackageWarning{korean-fonts}{Noto Serif CJK KR not found; Hangul Syllables will not render.}%
}

% Hangul Jamo (1100–11FF)
\IfFontExistsTF{Noto Serif CJK KR}{%
  \setTransitionsFor{HangulJamo}{\begingroup\hangulfont}{\endgroup}%
}{}

% Hangul Compatibility Jamo (3130–318F)
\IfFontExistsTF{Noto Serif CJK KR}{%
  \setTransitionsFor{HangulCompatibilityJamo}{\begingroup\hangulfont}{\endgroup}%
}{}

% Fixed-width Hangul for code blocks, if ever needed inline.
\IfFontExistsTF{Noto Sans Mono CJK KR}{%
  \newfontfamily\hangulmonofont{Noto Sans Mono CJK KR}%
}{%
  \PackageWarning{korean-fonts}{Noto Sans Mono CJK KR not found; falling back to \string\hangulfont\space for monospace Hangul.}%
}
